\documentclass[12pt]{article}
\usepackage[utf8]{inputenc}
\usepackage[T1]{fontenc}
\usepackage{hyperref}
\title{Fuliiru Language Resources for K-12 Curriculum Development}
\author{Compiled by Claude Sonnet 5.5 (claude-sonnet-5-5) for Andr\'as Kornai}
\date{}
\begin{document}
\maketitle

\section{Introduction}
Fuliiru (Kifuliiru; also Fuliiro, Fuliru, Fulero) is a Bantu language of South Kivu, Democratic Republic of the Congo, with about 400,000 speakers. Search results show a full Bible and some scripture media, but no school materials.

\section{Primary Educational Resources}
No primers or textbooks were found in web searches.

\section{Linguistic and Reference Materials}
\begin{itemize}
\item SemDom bibliographic record for a semantic-domain list based on Roger Van Otterloo's Kifuliiru dictionary (2000); described as unpublished, and no downloadable dictionary was found: \url{https://semdom.org/printpdf/16995}
\item Fuliru people, general background: \url{https://en.wikipedia.org/wiki/Fuliru_people}
\end{itemize}

\section{Religious and Community Resources}
\begin{itemize}
\item Full Bible, \emph{Ibibiriya: Igambo lya Rurema mu ndeto ye'kifuliiru}, on eBible.org: \url{https://ebible.org/study/content/texts/flr/about.html}
\item Scripture Earth index (Gospel of Mark, JESUS Film, Global Recordings audio, printed New Testament \emph{Ikihango Kihyahya}): \url{https://scriptureearth.org/00i-Scripture_Index.php?iso=flr}
\item Joshua Project people-group page: \url{https://joshuaproject.net/people_groups/11776}
\end{itemize}

\section{Digital and Audio Resources}
Audio and video scripture are listed on the Scripture Earth page above.

\section{Curriculum Development Context}
The language is not known to be taught in schools. The Bible text could supply a reading corpus and the Van Otterloo dictionary a lexical base if its owners release it; otherwise a curriculum would require new materials.
\end{document}